% Aportación separada, 9 de septiembre de 2026.
% RESULTADO_RECUPERADO: construcción sectorial del propietario 08_ckm_cp.tex.
% FORMALIZACION_REUNIDA: descomposición de paridad y compatibilidad representacional.
% Insertar antes de la carta racional CKM; conservar el desarrollo posterior.
% La procedencia y el alcance exacto están en 08b_procedencia_angular.md.

\subsection{Incidencia sectorial y construcción del lector angular}
\label{sec:ii09-construccion-sectorial}

La construcción angular utiliza las representaciones de acción y cierre
obtenidas previamente desde APP--TRIT--TPK. El ángulo directo, el
contraángulo y el defecto torsional conservan sus dominios propios.
El lector quark combina estas tres coordenadas mediante la incidencia
sectorial; su evaluación precede a las amplitudes complejas y al
contraste metrológico. La inclusión hexádico--octádica interviene en
el tercer sector a través de la diferencia de sus cardinales.

% ARQUITECTURA_AUTORAL_PREEXISTENTE: semicuadrado bidireccional de N42.
% FORMALIZACION_NUEVA: grupoide de pares y ponderacion por estabilizadores.
\subsubsection{Incidencia marcada y normalización bidireccional}
\label{sec:ii09-normalizacion-bidireccional}

Los invariantes del lector conservan la naturaleza del objeto del que
proceden. La intersección marcada $B_K=H_e\cap P_\pi$ tiene cardinal
$b=4$; el entero $q=7$ es la coordenada del pivote $p_\varphi$ en la
numeración orientada. El orden de determinación del diseño de Steiner
es $p=5$. Los soportes hexádico y octádico tienen cardinales $h_6=6$
y $o_8=8$, y el alfabeto de orientación trítica tiene cardinal $t=3$.
Un cardinal de soporte y una coordenada de pivote se utilizan mediante
sus aplicaciones respectivas, aun cuando ambos se publiquen como enteros.

La normalización del autocruce pentagonal admite la siguiente
realización explícita. Sea $P$ el soporte de cinco posiciones del cierre
y sea $X=P\times P$ su espacio de pares ordenados. La inversión
bidireccional sobre este espacio es
\[
 \jmath:X\longrightarrow X,\qquad \jmath(x,y)=(y,x),
 \qquad\jmath^2=I.
\]
Se conserva la información de estabilizador de cada órbita. Para una
órbita $[z]$ se define su peso como
$1/|\operatorname{Stab}_{C_2}(z)|$. La suma de esos pesos es la
cardinalidad del grupoide de acción $X/\!/C_2$.

\begin{proposition}[Medida orbital del autocruce bidireccional]
\label{prop:ii09-semicuadrado-orbital}
Para un soporte finito de cardinal $p$, la inversión de pares ordenados
determina
\[
 |(P\times P)/\!/C_2|
 =\sum_{[z]\in(P\times P)/C_2}
      \frac1{|\operatorname{Stab}_{C_2}(z)|}
 =\frac{p^2}{2}.
\]
Para $p=5$, el autocruce tiene diez órbitas libres y cinco órbitas
fijas, y su medida orbital es $10+5/2=25/2$.
\end{proposition}
\begin{proof}
Cada par diagonal $(x,x)$ forma una órbita fija, cuyo estabilizador
tiene orden dos; hay $p$ de ellas. Los $p(p-1)$ pares restantes se
agrupan de dos en dos, con estabilizador trivial. Su contribución es
por tanto $p(p-1)/2$, y la de la diagonal es $p/2$. La suma es
$p^2/2$. Equivalentemente, la identidad órbita--estabilizador da
$1/|\operatorname{Stab}(z)|=|[z]|/2$; sumar sobre todas las órbitas
produce $|P\times P|/2$.
\end{proof}

La medida conserva explícitamente la contribución de los puntos fijos.
El conjunto ordinario de órbitas tiene cardinal $p(p+1)/2$, que vale
quince en el cierre pentagonal. La medida orbital y ese cardinal son
dos lecturas determinadas de la misma inversión. La normalización
bidireccional del lector utiliza la primera: las diez órbitas libres
pesan una unidad y las cinco diagonales pesan media unidad.

Esta realización precisa el factor $p^2/2$ empleado en el sistema
sectorial siguiente. La asignación de esa contribución al sector $23$,
su orientación negativa y las otras dos reglas de incidencia permanecen
especificadas por el sistema sectorial completo. La prueba establece la
normalización del autocruce y preserva diferenciada la selección de los
funcionales angulares.


% FORMALIZACION_NUEVA de la elevacion de N39, compatible con (rho9,q9),
% T_d y Upd_t; arquitectura de residuo/cociente y orientacion preexistente.
% Fragmento de investigacion separado. No selecciona las reglas CKM N42.
\subsubsection{Transporte de acarreo y holonomía entera del cruce aritmético}
\label{subsec:transporte-entero-cruce}

La lectura modular del cruce admite una realización integral que mantiene
la información acumulada al atravesar el borde. Se desarrolla aquí esa
realización mediante los cocientes de transporte. Su dominio es una coordenatización
elevada del soporte APP; la posterior asignación de una región a un sector
angular conserva su propio mapa de incidencia.

\paragraph{Coordenatización nonádica y transporte elevado.}
Se emplea la descomposición exacta
\[
 \rho_9(n)=1+((n-1)\bmod9),\qquad
 q_9(n)=\left\lfloor\frac{n-1}{9}\right\rfloor,
 \qquad n=\rho_9(n)+9q_9(n).
\]
Una posición elevada es
$z=(x,y,q_x,q_y)\in D_9^2\times\mathbb Z^2$, donde
$X=x+9q_x$ e $Y=y+9q_y$. El incremento entero $u$ de la primera
coordenada actúa por
\[
 T^x_u(x,y,q_x,q_y)=
 \bigl(\rho_9(x+u),y,q_x+c_9(x;u),q_y\bigr),\qquad
 c_9(x;u)=\left\lfloor\frac{x+u-1}{9}\right\rfloor.
\]
La segunda coordenada se transporta de la misma manera. Los avances
cardinales corresponden a $u=\pm1$; el incremento nulo conserva el estado.
El cociclo satisface
\[
 c_9(x;u+v)=c_9(x;u)+c_9(\rho_9(x+u);v).
\]
En efecto, ambas partes son el cociente de la única descomposición de
$x+u+v$ con residuo en $D_9$. Por consiguiente,
$T^x_vT^x_u=T^x_{u+v}$ y $T^x_{-u}=(T^x_u)^{-1}$; los transportes
de las dos coordenadas conmutan. El estado conserva el número de cruces
de borde, incluso cuando la posición visible retorna.

\paragraph{Hojas elevadas y actualización de los registros.}
La realización integral se fija prolongando las operaciones aritméticas
APP a las coordenadas transportadas:
\begin{align*}
 \widetilde S(z)&=X+Y=x+y+9(q_x+q_y),\\
 \widetilde P(z)&=XY
 =xy+9(q_xy+q_yx)+81q_xq_y.
\end{align*}
En cada posición se conserva el registro de ambas hojas
\[
 \bigl(\rho_9(\widetilde S),q_9(\widetilde S),
       \rho_9(\widetilde P),q_9(\widetilde P)\bigr).
\]
Si $e:z\to z'$ es un enlace y $T$ la hoja seleccionada, su incremento
íntegro se obtiene del registro antes y después del transporte:
\begin{equation}
 \delta_e\widetilde T=
 \rho_9(\widetilde T(z'))-\rho_9(\widetilde T(z))
 +9\bigl(q_9(\widetilde T(z'))-q_9(\widetilde T(z))\bigr).
 \label{eq:cruce-incremento-registro}
\end{equation}
Esta identidad explicita la conservación: el salto del residuo y el
incremento del cociente reconstruyen juntos la variación de la hoja.
La selección de hoja, el transporte cardinal y la actualización del
registro son operaciones distintas. En esta coordenatización se compone primero la
selección ordenada de hojas $(T_x,T_y)$ con $T_d$ y con la actualización
de sus cocientes; la fórmula no postula que todo el TPK quede reducido
a esta representación local.

\begin{proposition}[Curvatura integral y transporte de frontera]
\label{prop:cruce-stokes-integral}
Defínanse
$\widetilde A_x=\delta_x\widetilde T_x$ y
$\widetilde A_y=\delta_y\widetilde T_y$, mediante
\eqref{eq:cruce-incremento-registro}, y sea
\[
 \widetilde F_{xy}(X,Y)=
 \widetilde A_x(X,Y)+\widetilde A_y(X+1,Y)
 -\widetilde A_x(X,Y+1)-\widetilde A_y(X,Y).
\]
Las elecciones homogéneas tienen curvatura cero. Las elecciones mixtas
$(\widetilde S,\widetilde P)$ y
$(\widetilde P,\widetilde S)$ tienen curvatura, respectivamente,
$+1$ y $-1$ como enteros. Para el borde positivo de un rectángulo
$R_{a,b}$ de lados enteros $a,b>0$,
\[
 \widetilde W_{SP}(\partial R_{a,b})=ab,\qquad
 \widetilde W_{PS}(\partial R_{a,b})=-ab.
\]
El resultado es independiente del punto base y del número de cruces
de las costuras nonádicas.
\end{proposition}

\begin{proof}
La reconstrucción exacta de ambas hojas da
\[
 \delta_x\widetilde S=\delta_y\widetilde S=1,
 \qquad\delta_x\widetilde P=Y,
 \qquad\delta_y\widetilde P=X.
\]
Por tanto, la conexión $SP$ es $(1,X)$, con curvatura
$1+(X+1)-1-X=1$; la conexión $PS$ es $(Y,1)$, con curvatura
$Y+1-(Y+1)-1=-1$. Para una hoja única se cancelan las diferencias
mixtas. La suma de curvaturas de las $ab$ plaquetas cancela cada
enlace interior con su inverso y conserva los enlaces de frontera.
Así se obtiene $\pm ab$ en $\mathbb Z$. Los cocientes transportados
garantizan que la misma identidad rige al atravesar cualquier costura.
\end{proof}

El acumulador de frontera puede conservarse a su vez en forma nonádica:
$W=r_W+9q_W$. Para una contribución entera $w_e$, se actualiza mediante
\[
 r_W'=\rho_9(r_W+w_e),\qquad
 q_W'=q_W+\left\lfloor\frac{r_W+w_e-1}{9}\right\rfloor.
\]
La composición por enlaces proporciona la misma suma que su evaluación
entera directa. En esta coordenatización, el cero entero tiene coordenadas
$(r_W,q_W)=(9,-1)$; esta es una representación aritmética del acumulador,
sin identificación adicional con los distintos ceros orientados del TRIT.

\paragraph{Tres operaciones sobre la orientación.}
Sea $\gamma$ una ruta cardinal cerrada. Su inversión temporal
$\gamma^{-1}$ invierte el orden de enlaces y la orientación de cada uno;
por aditividad,
\[
 \widetilde W_{SP}(\gamma^{-1})=-\widetilde W_{SP}(\gamma).
\]
El intercambio de las hojas satisface asimismo
\[
 \widetilde W_{PS}(\gamma)=-\widetilde W_{SP}(\gamma).
\]
Para comprobarlo en cualquier ruta, obsérvese que la suma de ambas
conexiones es el incremento exacto de
$\widetilde S+\widetilde P=X+Y+XY$; su integral sobre un lazo es cero.

La semivuelta simultánea de las dos direcciones actúa, por otra parte,
mediante $J_c(X,Y)=(2c_x-X,2c_y-Y)$, con $2c\in\mathbb Z^2$,
de modo que preserva la coordenatización reticular. Para un lazo se cumple
\[
 \widetilde W_{SP}(J_c\gamma)=\widetilde W_{SP}(\gamma).
\]
En efecto, la contribución horizontal cerrada es cero y la vertical
se transforma como
$\sum (2c_x-X)(-\delta Y)=\sum X\delta Y$, pues
$\sum\delta Y=0$. Una semivuelta espacial preserva, por tanto, la
orientación areal de esta coordenatización. Su combinación con la inversión temporal
o con el intercambio de hojas se calcula componiendo las operaciones
correspondientes.

La diferencia es pertinente para el registro electrónico: la inversión
simultánea de los dos desplazamientos cada $27$ avances produce, en
ventanas de $9$, la sucesión
\[
 \tau_j=(-1)^{\lfloor j/3\rfloor},\qquad 0\leq j<12,
 \qquad (+,+,+,-,-,-,+,+,+,-,-,-).
\]
Esta sucesión registra orientaciones de avance. La holonomía requiere
además los enlaces recorridos, las hojas seleccionadas y los cocientes
acumulados. Así, el transporte de la orientación electrónica se conserva
sin identificar una semivuelta de ambos ejes con la inversión temporal
de un lazo.

\paragraph{Aplicación al cruce marcado y alcance sectorial.}
Para el cruce rectangular de lados $4$ y $7$, esta realización evalúa
exactamente
\[
 \widetilde W_{SP}=28=1+9\cdot3,\qquad
 \widetilde W_{SP}(\gamma^{-1})=-28=8+9(-4).
\]
El entero conserva el área orientada; el residuo aislado representa sólo
su clase nonádica. Por ejemplo, desde $(7,1)$ los representantes
periódicos de los enlaces suman $-35$, mientras que la corrección de
acarreo añade $63$ y recupera $28$.

En la incidencia marcada del lector angular, el $4$ es
$|H_e\cap P_\pi|$ y el $7$ es el pivote marcado $p_\varphi$.
La lectura rectangular aporta un realizador entero del producto
$|H_e\cap P_\pi|\,p_\varphi$ una vez declarado ese cruce. El
mapa que identifique esa región con el sector $12$ debe conservar
la marca y la orientación del sistema sectorial de
la Definición~\ref{def:ii09-sistema-sectorial}. Las asignaciones
sectoriales $12,23,13$ y sus combinaciones de $A,h,\Delta^\circ$
son una especificación adicional a la ley de holonomía rectangular.
El presente resultado materializa la elevación de área, el acarreo y
sus involuciones, con esos datos de partida explícitos.


\begin{definition}[Sistema sectorial de incidencia angular]
\label{def:ii09-sistema-sectorial}
El sistema sectorial consta del sexteto ordenado de invariantes
\[
 (b,p,q,h_6,o_8,t)=(4,5,7,6,8,3)
\]
y de las aplicaciones lineales, sobre coordenadas angulares
\((a,u,d)\),
\begin{align}
 \mathscr A_{12}^{\rm ang}(a,u,d)&=2a-pu+bq\,d,\label{eq:ii09-sector12}\\
 \mathscr A_{23}^{\rm ang}(a,u,d)&=qu-\frac{p^2}{2}\,d,\label{eq:ii09-sector23}\\
 \mathscr A_{13}^{\rm ang}(a,u,d)&=a-pb\,u-\frac{o_8-h_6}{t}\,d.
 \label{eq:ii09-sector13}
\end{align}
Los símbolos \(h_6,o_8\) conservan los cardinales de la incidencia
marcada; \(t=3\) es la normalización trítica declarada de ese sector.
El sexteto y las tres reglas forman conjuntamente el sistema de
partida de esta construcción.
\end{definition}

\begin{theorem}[Evaluación exacta del sistema sectorial]
\label{thm:ii09-matriz-sectorial}
La aplicación conjunta
\((\mathscr A_{12}^{\rm ang},\mathscr A_{23}^{\rm ang},
\mathscr A_{13}^{\rm ang})^{\mathsf T}\) de la
Definición~\ref{def:ii09-sistema-sectorial} tiene por matriz única
\[
 \mathsf M_{\rm CKM}=
 \begin{pmatrix}
 2&-5&28\\
 0&7&-25/2\\
 1&-20&-2/3
 \end{pmatrix}.
\]
Las tres columnas son las imágenes de los vectores coordenados
\((1,0,0)\), \((0,1,0)\) y \((0,0,1)\) bajo esas aplicaciones.
\end{theorem}
\begin{proof}
Las operaciones sectoriales dan
\[
 bq=4\cdot7=28,\qquad p^2/2=5^2/2=25/2,
 \qquad pb=5\cdot4=20,\qquad
 (o_8-h_6)/t=(8-6)/3=2/3.
\]
Por tanto, las imágenes sucesivas de la base coordenada son
\[
 c_A=(2,0,1)^{\mathsf T},\qquad
 c_h=(-5,7,-20)^{\mathsf T},\qquad
 c_\Delta=(28,-25/2,-2/3)^{\mathsf T}.
\]
Una aplicación lineal queda determinada por sus imágenes sobre una
base. La matriz cuyas columnas son estas imágenes es la indicada.
La prueba utiliza el sistema sectorial completo de la definición.
\end{proof}

Sobre las coordenadas previamente producidas se toma
\[
 (a,u,d)=\left(A_{\deg},\frac{C^*_{\deg}}6,
                       \frac{180}{\pi}\Delta_4\right).
\]
Así quedan individualizadas las contribuciones directa, conjugada
y torsional a cada ángulo. En particular, la diferencia
\(o_8-h_6=2\) se conserva en el coeficiente del defecto torsional
del sector \(13\). Esta presencia concreta permite seguir la
incidencia hasta el lector angular, manteniendo separados el
sistema sectorial, su evaluación y la parametrización unitaria.
